%% ODER: format ==         = "\mathrel{==}"
%% ODER: format /=         = "\neq "
%
%
\makeatletter
\@ifundefined{lhs2tex.lhs2tex.sty.read}%
  {\@namedef{lhs2tex.lhs2tex.sty.read}{}%
   \newcommand\SkipToFmtEnd{}%
   \newcommand\EndFmtInput{}%
   \long\def\SkipToFmtEnd#1\EndFmtInput{}%
  }\SkipToFmtEnd

\newcommand\ReadOnlyOnce[1]{\@ifundefined{#1}{\@namedef{#1}{}}\SkipToFmtEnd}
\usepackage{amstext}
\usepackage{amssymb}
\usepackage{stmaryrd}
\DeclareFontFamily{OT1}{cmtex}{}
\DeclareFontShape{OT1}{cmtex}{m}{n}
  {<5><6><7><8>cmtex8
   <9>cmtex9
   <10><10.95><12><14.4><17.28><20.74><24.88>cmtex10}{}
\DeclareFontShape{OT1}{cmtex}{m}{it}
  {<-> ssub * cmtt/m/it}{}
\newcommand{\texfamily}{\fontfamily{cmtex}\selectfont}
\DeclareFontShape{OT1}{cmtt}{bx}{n}
  {<5><6><7><8>cmtt8
   <9>cmbtt9
   <10><10.95><12><14.4><17.28><20.74><24.88>cmbtt10}{}
\DeclareFontShape{OT1}{cmtex}{bx}{n}
  {<-> ssub * cmtt/bx/n}{}
\newcommand{\tex}[1]{\text{\texfamily#1}}	% NEU

\newcommand{\Sp}{\hskip.33334em\relax}


\newcommand{\Conid}[1]{\mathit{#1}}
\newcommand{\Varid}[1]{\mathit{#1}}
\newcommand{\anonymous}{\kern0.06em \vbox{\hrule\@width.5em}}
\newcommand{\plus}{\mathbin{+\!\!\!+}}
\newcommand{\bind}{\mathbin{>\!\!\!>\mkern-6.7mu=}}
\newcommand{\rbind}{\mathbin{=\mkern-6.7mu<\!\!\!<}}% suggested by Neil Mitchell
\newcommand{\sequ}{\mathbin{>\!\!\!>}}
\renewcommand{\leq}{\leqslant}
\renewcommand{\geq}{\geqslant}
\usepackage{polytable}

%mathindent has to be defined
\@ifundefined{mathindent}%
  {\newdimen\mathindent\mathindent\leftmargini}%
  {}%

\def\resethooks{%
  \global\let\SaveRestoreHook\empty
  \global\let\ColumnHook\empty}
\newcommand*{\savecolumns}[1][default]%
  {\g@addto@macro\SaveRestoreHook{\savecolumns[#1]}}
\newcommand*{\restorecolumns}[1][default]%
  {\g@addto@macro\SaveRestoreHook{\restorecolumns[#1]}}
\newcommand*{\aligncolumn}[2]%
  {\g@addto@macro\ColumnHook{\column{#1}{#2}}}

\resethooks

\newcommand{\onelinecommentchars}{\quad-{}- }
\newcommand{\commentbeginchars}{\enskip\{-}
\newcommand{\commentendchars}{-\}\enskip}

\newcommand{\visiblecomments}{%
  \let\onelinecomment=\onelinecommentchars
  \let\commentbegin=\commentbeginchars
  \let\commentend=\commentendchars}

\newcommand{\invisiblecomments}{%
  \let\onelinecomment=\empty
  \let\commentbegin=\empty
  \let\commentend=\empty}

\visiblecomments

\newlength{\blanklineskip}
\setlength{\blanklineskip}{0.66084ex}

\newcommand{\hsindent}[1]{\quad}% default is fixed indentation
\let\hspre\empty
\let\hspost\empty
\newcommand{\NB}{\textbf{NB}}
\newcommand{\Todo}[1]{$\langle$\textbf{To do:}~#1$\rangle$}

\EndFmtInput
\makeatother
%
%
%
%
%
%
% This package provides two environments suitable to take the place
% of hscode, called "plainhscode" and "arrayhscode". 
%
% The plain environment surrounds each code block by vertical space,
% and it uses \abovedisplayskip and \belowdisplayskip to get spacing
% similar to formulas. Note that if these dimensions are changed,
% the spacing around displayed math formulas changes as well.
% All code is indented using \leftskip.
%
% Changed 19.08.2004 to reflect changes in colorcode. Should work with
% CodeGroup.sty.
%
\ReadOnlyOnce{polycode.fmt}%
\makeatletter

\newcommand{\hsnewpar}[1]%
  {{\parskip=0pt\parindent=0pt\par\vskip #1\noindent}}

% can be used, for instance, to redefine the code size, by setting the
% command to \small or something alike
\newcommand{\hscodestyle}{}

% The command \sethscode can be used to switch the code formatting
% behaviour by mapping the hscode environment in the subst directive
% to a new LaTeX environment.

\newcommand{\sethscode}[1]%
  {\expandafter\let\expandafter\hscode\csname #1\endcsname
   \expandafter\let\expandafter\endhscode\csname end#1\endcsname}

% "compatibility" mode restores the non-polycode.fmt layout.

\newenvironment{compathscode}%
  {\par\noindent
   \advance\leftskip\mathindent
   \hscodestyle
   \let\\=\@normalcr
   \let\hspre\(\let\hspost\)%
   \pboxed}%
  {\endpboxed\)%
   \par\noindent
   \ignorespacesafterend}

\newcommand{\compaths}{\sethscode{compathscode}}

% "plain" mode is the proposed default.
% It should now work with \centering.
% This required some changes. The old version
% is still available for reference as oldplainhscode.

\newenvironment{plainhscode}%
  {\hsnewpar\abovedisplayskip
   \advance\leftskip\mathindent
   \hscodestyle
   \let\hspre\(\let\hspost\)%
   \pboxed}%
  {\endpboxed%
   \hsnewpar\belowdisplayskip
   \ignorespacesafterend}

\newenvironment{oldplainhscode}%
  {\hsnewpar\abovedisplayskip
   \advance\leftskip\mathindent
   \hscodestyle
   \let\\=\@normalcr
   \(\pboxed}%
  {\endpboxed\)%
   \hsnewpar\belowdisplayskip
   \ignorespacesafterend}

% Here, we make plainhscode the default environment.

\newcommand{\plainhs}{\sethscode{plainhscode}}
\newcommand{\oldplainhs}{\sethscode{oldplainhscode}}
\plainhs

% The arrayhscode is like plain, but makes use of polytable's
% parray environment which disallows page breaks in code blocks.

\newenvironment{arrayhscode}%
  {\hsnewpar\abovedisplayskip
   \advance\leftskip\mathindent
   \hscodestyle
   \let\\=\@normalcr
   \(\parray}%
  {\endparray\)%
   \hsnewpar\belowdisplayskip
   \ignorespacesafterend}

\newcommand{\arrayhs}{\sethscode{arrayhscode}}

% The mathhscode environment also makes use of polytable's parray 
% environment. It is supposed to be used only inside math mode 
% (I used it to typeset the type rules in my thesis).

\newenvironment{mathhscode}%
  {\parray}{\endparray}

\newcommand{\mathhs}{\sethscode{mathhscode}}

% texths is similar to mathhs, but works in text mode.

\newenvironment{texthscode}%
  {\(\parray}{\endparray\)}

\newcommand{\texths}{\sethscode{texthscode}}

% The framed environment places code in a framed box.

\def\codeframewidth{\arrayrulewidth}
\RequirePackage{calc}

\newenvironment{framedhscode}%
  {\parskip=\abovedisplayskip\par\noindent
   \hscodestyle
   \arrayrulewidth=\codeframewidth
   \tabular{@{}|p{\linewidth-2\arraycolsep-2\arrayrulewidth-2pt}|@{}}%
   \hline\framedhslinecorrect\\{-1.5ex}%
   \let\endoflinesave=\\
   \let\\=\@normalcr
   \(\pboxed}%
  {\endpboxed\)%
   \framedhslinecorrect\endoflinesave{.5ex}\hline
   \endtabular
   \parskip=\belowdisplayskip\par\noindent
   \ignorespacesafterend}

\newcommand{\framedhslinecorrect}[2]%
  {#1[#2]}

\newcommand{\framedhs}{\sethscode{framedhscode}}

% The inlinehscode environment is an experimental environment
% that can be used to typeset displayed code inline.

\newenvironment{inlinehscode}%
  {\(\def\column##1##2{}%
   \let\>\undefined\let\<\undefined\let\\\undefined
   \newcommand\>[1][]{}\newcommand\<[1][]{}\newcommand\\[1][]{}%
   \def\fromto##1##2##3{##3}%
   \def\nextline{}}{\) }%

\newcommand{\inlinehs}{\sethscode{inlinehscode}}

% The joincode environment is a separate environment that
% can be used to surround and thereby connect multiple code
% blocks.

\newenvironment{joincode}%
  {\let\orighscode=\hscode
   \let\origendhscode=\endhscode
   \def\endhscode{\def\hscode{\endgroup\def\@currenvir{hscode}\\}\begingroup}
   %\let\SaveRestoreHook=\empty
   %\let\ColumnHook=\empty
   %\let\resethooks=\empty
   \orighscode\def\hscode{\endgroup\def\@currenvir{hscode}}}%
  {\origendhscode
   \global\let\hscode=\orighscode
   \global\let\endhscode=\origendhscode}%

\makeatother
\EndFmtInput
%



%% ----------- Applications --------------
%% %format applyx a b = a "@" b
%% %format applyx a b = a "\," b

%% Function application with no arguments

%% Function application with 1 argument

%% Function application with 1 argument, without parenthesis

%% Function application with many arguments

%% ----------- End of Applications --------------


%% Sequences e1; ...; en; v

%% xs etc stole syntax used in examples.

%% Old version: | for BNF, [] for choice
%%%format choice = "[ \! ]"
%%%format `choice` = "\mathop{[ \! ]}"

%% New version: tall | for BNF, fat | for choice
%%format vbar = "\mathop{\bigm|}"
%%format choice = "\boldsymbol{|}"
%%format `choice` = "\mathop{\boldsymbol{|}}"




%% Arrays and tuples
%% %format array (x) = "\textbf{array}\{" x "\}"
%% %format arrKW = "\textbf{arr}"
%% %format arr (x) = arrKW "\{" x "\}"

%% %format defKW = "\textbf{def}"

% only used in comments, but still prettier this way:
%%%%format map = "\textbf{map}"
% format := = "\mapsto"
%  format effs x = "\!\!<\!\!" x "\!\!>\!\!" dots































%% ----------- Metatheory --------------
%%%%%%%%%%%%%%%%%%%%%%%%%%%%%%%%%%%%

% Typesetting rewrite rules
\newcommand{\rulename}[1]{\textsc{\small #1}}
\newcommand{\vb}{\; \ensuremath{\mathop{\mid}} \;}    % Vertical bar
\newcommand{\hole}{\square}
\newcommand{\freevars}[1]{\mathsf{fvs}(#1)}
\newcommand{\boundvars}[1]{\mathsf{bvs}(#1)}
\newcommand{\valvars}[1]{\mathsf{vvs}(#1)}
\newcommand{\disjoint}{\;\#\;}
\newcommand{\hnf}{\mathit{hnf}}  % head-normal form
\newcommand{\op}{\mathit{op}}
\newcommand{\subst}[3]{#1\{#3/#2\}}

% Name of the language
\newcommand{\versecalc}{$\mathcal{VC}$}
\newcommand{\calcname}{\versecalc\xspace}
\newcommand{\mypara}[1]{\smallskip\noindent\textbf{#1}\xspace}

% \newcommand{\calcname}{Fresh Calculus\xspace}
% \newcommand{\lcalcname}{Lenient Calculus\xspace}
% \newcommand{\calcshort}{FC\xspace}
% \newcommand{\vcname}{Verse Calculus\xspace}
% \newcommand{\vcshort}{VC\xspace}

% Right arrow for rewrites
\newcommand{\movesto}{\longrightarrow}
\newcommand{\movestoX}{\rightsquigarrow}
\newcommand{\movesfromX}{\leftsquigarrow}
\newcommand{\movesfrom}{\longleftarrow}
\newcommand{\convertible}{\longleftrightarrow^{\ast}}

% Arrows with explanation, including the common cases of one, two or a repeated rule
\newcommand{\movestoC}[1]{\movesto\!\!\{#1\}}
\newcommand{\movestoR}[1]{\movestoC{\rulename{#1}}}
\newcommand{\movestoRR}[2]{\movestoC{\rulename{#1}, \rulename{#2}}}
\newcommand{\movestoRn}[2]{\movestoC{\rulename{#1} \times #2}}


\newcommand{\rightarrowdbl}{\rightarrow\mathrel{\mkern-14mu}\rightarrow}
\newcommand{\xrightarrowdbl}[2][]{%
  \xrightarrow[#1]{#2}\mathrel{\mkern-14mu}\rightarrow
}

\newcommand{\rr}{\ensuremath{\mathcal{R}}\xspace}
\newcommand{\rs}{\ensuremath{\mathcal{S}}\xspace}
\newcommand{\goone}[1]{\xrightarrow{\epsilon}_{#1}}
\newcommand{\go}[1]{\xrightarrow{}_{#1}}
\newcommand{\gos}[1]{\xrightarrowdbl{}_{#1}}
\newcommand{\gok}[2]{\xrightarrowdbl{#2}_{#1}}
\newcommand{\rstepo}[3]{\ensuremath{#2 \goone{#1} #3}}
\newcommand{\rstep}[3]{\ensuremath{#2 \go{#1} #3}}
\newcommand{\rsteps}[3]{\ensuremath{#2 \gos{#1} #3}}
\newcommand{\rstepk}[4]{\ensuremath{#3 \gok{#1}{#2} #4}}
\newcommand{\univ}{\ensuremath{\mathcal{E}}}
\newcommand{\join}[3]{{#2} \sim_{#1} {#3}}
\newcommand{\dia}{\ensuremath{\diamond}}

\newcommand{\rnk}{\ensuremath{\rho}}
\newcommand{\rord}{\ensuremath{\mathcal{O}}}
\newcommand{\runi}{\ensuremath{\mathcal{U}}}
\newcommand{\rapp}{\ensuremath{\mathcal{A}}}
\newcommand{\rcon}{\ensuremath{\mathcal{N}}}
\newcommand{\rgc}{\ensuremath{\mathcal{G}}}
\newcommand{\rfail}{\ensuremath{\mathcal{F}}}
\newcommand{\rspc}{\ensuremath{\mathcal{S}}}

\newcommand{\rdx}[1]{\ensuremath{\Delta_{#1}}}
\newcommand{\tx}[1]{\ensuremath{\rstep{#1}{\rdx{#1}}{\rdx{#1}'}}}
\newcommand{\txi}[1]{\ensuremath{\rstep{}{\rdx{#1}}{\rdx{#1}'}}}


\newcommand{\wf}[1]{\ensuremath{\vdash #1}}
\newcommand{\wff}[5]{\ensuremath{#1 \vdash #2 \rightsquigarrow (#3 \mid #4 \mid #5)}}
\newcommand{\wffv}[4]{\ensuremath{\vdash_\ensuremath{\Varid{v}} #1 \rightsquigarrow (#2 \mid #3 \mid #4)}}
\newcommand{\wffmany}[4]{\ensuremath{\vdash_{\!*} #1 \rightsquigarrow (#2 \mid #3 \mid #4)}}
% \newcommand{\lvc}[1]{\colorbox{lightgray}{#1}}
% \newcommand{\lvc}[1]{\colorbox{lightgray}{#1}}
\newcommand{\lvc}[1]{#1}
\newcommand{\relDescription}[1]{\ensuremath{\textrm{\textbf{#1}}}}
\newcommand{\judgementHead}[2]{\ensuremath{\relDescription{#1}\hfill\fbox{#2}}}
\newcommand{\judgementHeadNameOnly}[1]{\ensuremath{\relDescription{#1}\hfill}}


\newcommand{\rulegroupname}[1]{\multicolumn{4}{l}{\textit{#1}}\\}
\newcommand{\largerhs}[1]{\multicolumn{2}{@{}l}{#1}}


% --------------------------------------------------------

\newcommand{\figureSyntax}{
\begin{figure}[t]

$$
\begin{array}{l}
  \textbf{Abstract syntax} \\
    \begin{array}{lrcl}
    \textit{Integers}      & k \\
    \textit{Variables}     & \multicolumn{3}{l}{x,y,z,f,g} \\
    \textit{Primops}       & \op   & ::= & \ensuremath{\textbf{gt}} \vb \ensuremath{\textbf{add}} \\
    \textit{Scalar Values} & s     & ::= & x \vb k \vb op \\
    \textit{Heap Values}   & h     & ::= & \ensuremath{\langle{}\Varid{s}_{\mathrm{1}},\cdots,\Varid{s}_{\Varid{n}}\rangle} \vb \ensuremath{\lambda\Varid{x}.\,\Varid{e}} \\
    \textit{Head Values}   & \ensuremath{\mathit{hnf}} & ::= & h \vb k \\
    \textit{Values}        & \ensuremath{\Varid{v}}   & ::= & s \vb h \\
    \textit{Expressions}   & e     & ::= & \ensuremath{\Varid{v}} \vb \ensuremath{eu;\,\Varid{e}} \vb \ensuremath{\exists\Varid{x}.\,\Varid{e}} \vb \ensuremath{\textbf{fail}}
                                             \vb \ensuremath{\Varid{e}_{\mathrm{1}}\;\rule[-0.3mm]{0.5mm}{2.5mm}\;\Varid{e}_{\mathrm{2}}} \vb \ensuremath{\Varid{v}_{\mathrm{1}}\,\Varid{v}_{\mathrm{2}}} \vb \ensuremath{\textbf{one}\{\Varid{e}\}} \vb \ensuremath{\textbf{all}\{\Varid{e}\}} \\
                           & eu     & ::= & \ensuremath{\Varid{e}} \vb \ensuremath{\Varid{v}=\Varid{e}} \\
%     \textit{Sequences}     & |se|  & ::= & v \vb |e ; se| \\
    \textit{Programs}      & p     & ::= & \ensuremath{\textbf{one}\{\Varid{e}\}} \quad \text{where} \; \freevars{e} = \emptyset \\
    \textit{Bindings}      & c     & ::= & \ensuremath{\Varid{x}\mathrel{=}\Varid{v}} \\
%     \multicolumn{4}{l}{\text{\lennart{Do we want to add |int| to the primops?  It's in the denotational semantics, and somewhat interesting.}}}
\end{array} \\
\\[-3mm]
\textbf{Concrete syntax:} \;\;
    \begin{array}[t]{l}
          \text{Infix operators ``\ensuremath{\rule[-0.3mm]{0.5mm}{2.5mm}}'', ``;'', ``\ensuremath{\mathrel{=}}'', and ``>'' are all right-associative.} \\
          \text{``\ensuremath{\mathrel{=}}'' binds more tightly than ``;''.} \\
          \text{Function application \ensuremath{(\Varid{v}_{\mathrm{1}}\;\Varid{v}_{\mathrm{2}})} is left-associatve, as usual.} \\
          \text{``$\lambda$'', ``\ensuremath{\exists}'' scope as far to the right as possible.} \\
          \hspace{2mm} \text{{\eg}, \ensuremath{(\lambda\Varid{y}.\,\exists\Varid{x}.\,\Varid{x}\mathrel{=}\mathrm{1};\,\Varid{x}\mathbin{+}\Varid{y})} means \ensuremath{(\lambda\Varid{y}.\,(\exists\Varid{x}.\,((\Varid{x}\mathrel{=}\mathrm{1});\,(\Varid{x}\mathbin{+}\Varid{y}))))}.}
    \end{array} \\
\textbf{Desugaring} \\
    \begin{array}{rcll}
%%      |emptytup| & \text{means} & |arr ()|\\
%%      |tup (v_1, dots, v_n)| & \text{means} & |arr (v_1,dots, v_n)| \quad n \geq 2 \\
      v_1 + v_2 & \text{means} & \ensuremath{\textbf{add}\langle\Varid{v}_{\mathrm{1}},\Varid{v}_{\mathrm{2}}\rangle} & \\
      v_1 > v_2 & \text{means} & \ensuremath{\textbf{gt}\langle\Varid{v}_{\mathrm{1}},\Varid{v}_{\mathrm{2}}\rangle} & \\
      \ensuremath{\exists\Varid{x}_{\mathrm{1}}\hspace{0.3mm}\Varid{x}_{\mathrm{2}}\;\cdots\;\Varid{x}_{\Varid{n}}.\,\Varid{e}} & \text{means} & \ensuremath{\exists\Varid{x}_{\mathrm{1}}.\,\exists\Varid{x}_{\mathrm{2}}.\,\cdots \exists\Varid{x}_{\Varid{n}}.\,\Varid{e}} & \\
      \ensuremath{\Varid{e}_{\mathrm{1}}(\Varid{e}_{\mathrm{2}})} & \text{means} & \ensuremath{\exists\Varid{f}\;\Varid{a}.\,\Varid{f}=\Varid{e}_{\mathrm{1}};\,\Varid{a}=\Varid{e}_{\mathrm{2}};\,\Varid{f}(\Varid{a})}  & \ensuremath{\Varid{f},\Varid{a}}\ \text{fresh} \\
      \ensuremath{\langle{}\Varid{e}_{\mathrm{1}},\cdots,\Varid{e}_{\Varid{n}}\rangle} & \text{means} &
          \ensuremath{\exists\Varid{x}_{\mathrm{1}}\hspace{0.3mm}\Varid{x}_{\mathrm{2}}\;\cdots\;\Varid{x}_{\Varid{n}}.\,\Varid{x}_{\mathrm{1}}\mathrel{=}\Varid{e}_{\mathrm{1}};\,\cdots;\,\Varid{x}_{\Varid{n}}\mathrel{=}\Varid{e}_{\Varid{n}};\,\langle{}\Varid{x}_{\mathrm{1}},\cdots,\Varid{x}_{\Varid{n}}\rangle} & \text{\ensuremath{\Varid{x}_{\Varid{i}}} fresh} \\
      \ensuremath{\Varid{e}_{\mathrm{1}}=\Varid{e}_{\mathrm{2}}}      & \text{means} & \ensuremath{\exists\Varid{x}.\,\Varid{x}=\Varid{e}_{\mathrm{1}};\,\Varid{x}=\Varid{e}_{\mathrm{2}};\,\Varid{x}}       & \ensuremath{\Varid{x}}\ \text{fresh} \\
%      |e_1 ; e_2|      & \text{means} & |def x (e_1; x == e_2 ; x)|       & |x|\ \text{fresh} \\

      \ensuremath{\textbf{if}\;\Varid{e}_{\mathrm{1}}\;\textbf{then}\;\Varid{e}_{\mathrm{2}}\;\textbf{else}\;\Varid{e}_{\mathrm{3}}} & \text{means}
      & \ensuremath{\exists\Varid{v}.\,\Varid{v}=\textbf{one}\{(\Varid{e}_{\mathrm{1}};\,\lambda x.\,\Varid{e}_{\mathrm{2}})\;\rule[-0.3mm]{0.5mm}{2.5mm}\;(\lambda x.\,\Varid{e}_{\mathrm{3}})\};\,\Varid{v}\langle\rangle} & \text{\ensuremath{\Varid{v},\Varid{x}} fresh} \\
%       |for^(e_1) do e_2| & \text{means} & |def v (v == all (e_1; thunk e_2); applyN map (lam z (apply0 z), v))|
      \ensuremath{\Varid{x}\mathrel{\coloneqq}\Varid{e}_{\mathrm{1}};\,\Varid{e}_{\mathrm{2}}} & \text{means} & \ensuremath{\exists\Varid{x}.\,\Varid{x}=\Varid{e}_{\mathrm{1}};\,\Varid{e}_{\mathrm{2}}} \\
     \end{array} \\
  \\[-2mm]
  \text{$\freevars{e}$ means the free variable of $e$; in \versecalc{}, $\lambda$ and \ensuremath{\exists} are the only binders.}
  %\text{$S_1 \disjoint S_2$ means that set $S_1$ is disjoint from $S_2$.}
  \end{array}
$$
\caption{\textbf{The Verse Calculus: Syntax}}
\label{fig:syntax}
\end{figure}
}

% --------------------------------------------------------

\newcommand{\figureRewriteRules}{
\begin{figure}
$$\begin{array}{lr@{\;}c@{\;}l}
\textit{Expression context}     & E  & ::= & \hole \vb \ensuremath{\langle{}\Varid{s}_{\mathrm{1}},\cdots,\hole,\cdots,\Varid{s}_{\Varid{n}}\rangle} \vb \ensuremath{\lambda\Varid{x}.\,\Conid{E}} \vb \ensuremath{\exists\Varid{x}.\,\Conid{E}} 
                                              \vb \ensuremath{\Conid{E}=\Varid{e}} \vb \ensuremath{\Varid{e}=\Conid{E}} \\
                                &    & \vb & \ensuremath{\Conid{E};\,\Varid{e}} \vb \ensuremath{\Varid{e};\,\Conid{E}} \vb \ensuremath{\Conid{E}\,\Varid{v}} \vb \ensuremath{\Varid{v}\,\Conid{E}}
                                       \vb \ensuremath{\Conid{E}\;\rule[-0.3mm]{0.5mm}{2.5mm}\;\Varid{e}} \vb \ensuremath{\Varid{e}\;\rule[-0.3mm]{0.5mm}{2.5mm}\;\Conid{E}} \vb \ensuremath{\textbf{all}\{\Conid{E}\}} \vb \ensuremath{\textbf{one}\{\Conid{E}\}} \\
\textit{Application context}    & A  & ::= & \ensuremath{\hole\,\Varid{v}} \vb \ensuremath{\Varid{op}\,\hole} \vb \ensuremath{\hole=\mathit{hnf}} \vb \ensuremath{\Varid{v}=\Conid{A}} \vb \ensuremath{\exists\Varid{x}.\,\Conid{A}} \vb \ensuremath{\Conid{A};\,\Varid{e}} \vb \ensuremath{\Varid{e};\,\Conid{A}} \\
                                &    & \vb & \ensuremath{\Conid{A}\;\rule[-0.3mm]{0.5mm}{2.5mm}\;\Varid{e}} \vb \ensuremath{\Varid{e}\;\rule[-0.3mm]{0.5mm}{2.5mm}\;\Conid{A}} \vb \ensuremath{\textbf{all}\{\Conid{A}\}} \vb \ensuremath{\textbf{one}\{\Conid{A}\}} \\
\textit{Scope context}          & SX & ::= & \ensuremath{\hole\;\rule[-0.3mm]{0.5mm}{2.5mm}\;\Varid{e}} \vb \ensuremath{\Varid{e}\;\rule[-0.3mm]{0.5mm}{2.5mm}\;\hole} \vb \ensuremath{\textbf{one}\{\hole\}} \vb \ensuremath{\textbf{all}\{\hole\}} \\
\textit{Choice context}         & CX & ::= & \hole \vb \ensuremath{\Varid{v}=CX} \vb \ensuremath{CX;\,\Varid{e}} \vb \ensuremath{\Varid{ce};\,CX}  \vb \ensuremath{\exists\Varid{x}.\,CX} \\
\textit{Choice-free expr}      & ce & ::= & \ensuremath{\Varid{v}} \vb \ensuremath{\Varid{v}=\Varid{ce}} \vb \ensuremath{\Varid{ce}_{\mathrm{1}};\,\Varid{ce}_{\mathrm{2}}} \vb \ensuremath{\textbf{one}\{\Varid{e}\}} \vb \ensuremath{\textbf{all}\{\Varid{e}\}}
                                             %%% NOTE: the following is not in the POPL submission,
                                             %%% but needed to make progress in (1,2)[x:int]
                                             \vb \ensuremath{\Varid{op}(\Varid{v})}
                                             \vb \ensuremath{\exists\Varid{x}.\,\Varid{ce}} \\
\textit{Bound variables} & \multicolumn{3}{l}{
  \boundvars{E} =
  \begin{array}[t]{l}
      \text{The variables that are bound by \ensuremath{\Conid{E}} at the hole} \\
      \text{e.g. $\boundvars{\ensuremath{(\exists\Varid{x}.\,\Varid{x}\mathrel{=}\mathrm{3})\hspace{1mm}\mathop{\rule[-0.3mm]{0.5mm}{2.5mm}}\hspace{1mm}(\exists\Varid{y}.\,\hole\mathrel{=}\mathrm{4}))}} = \{y\}$}
  \end{array}
   }
\end{array}$$
\vspace{-6mm}
\begin{longtable}[l]{@{\quad} l @{\hspace{-.5em}} r >{$}c<{$} l l}
\rulegroupname{Unification:~\runi}
\rulename{deref-s}  & \ensuremath{\Varid{x}=\Varid{s};\,\Conid{E} [\mskip1.5mu \Varid{x}\mskip1.5mu]} & \movesto & \ensuremath{\Varid{x}=\Varid{s};\,\Conid{E} [\mskip1.5mu \Varid{s}\mskip1.5mu]} & $x \not \equiv s, x \not\in \boundvars{\ensuremath{\Conid{E}}}, s \not\in \boundvars{\ensuremath{\Conid{E}}}$ \\
\rulename{deref-h}  & \ensuremath{\Varid{x}=\Varid{h};\,\Conid{A} [\mskip1.5mu \Varid{x}\mskip1.5mu]} & \movesto & \ensuremath{\Varid{x}=\Varid{h};\,\Conid{A} [\mskip1.5mu \Varid{h}\mskip1.5mu]} & $x \not\in \boundvars{A}, \freevars{h} \not\in \boundvars{A}$ \\
\rulename{u-scalar} & \ensuremath{\Varid{s}=\Varid{s};\,\Varid{e}}                                    & \movesto & e \\
\rulename{u-tup}    & \ensuremath{\langle{}\Varid{v}_{\mathrm{1}},\cdots,\Varid{v}_{\Varid{n}}\rangle=\langle{}\Varid{v}_{1}',\cdots,\Varid{v}_{\Varid{n}}'\rangle;\,\Varid{e}} & \movesto & \multicolumn{2}{l} {\ensuremath{\Varid{v}_{\mathrm{1}}=\Varid{v}_{1}';\,\cdots;\,\Varid{v}_{\Varid{n}}=\Varid{v}_{\Varid{n}}';\,\Varid{e}}} \\
\rulename{u-fail}   & \ensuremath{\mathit{hnf}_{\mathrm{1}}=\mathit{hnf}_{\mathrm{2}}} & \movesto & \ensuremath{\textbf{fail}} & \text{if neither \rulename{u-scalar} nor \rulename{u-tup} match} \\
\end{longtable}
\vspace{-4mm}
\begin{longtable}[l]{@{\quad} l r >{$}c<{$} l l}
\rulegroupname{Application:~\rapp}
\rulename{app-beta} & \ensuremath{(\lambda\Varid{x}.\,\Varid{e})\,\Varid{v}} & \movesto & \ensuremath{\exists\Varid{x}.\,\Varid{x}\mathrel{=}\Varid{v};\,\Varid{e}}  & if $\ensuremath{\Varid{x}} \not \in \freevars{\ensuremath{\Varid{v}}}$ \\
\rulename{app-tup0} & \ensuremath{\langle{}\rangle\,\Varid{v}} & \movesto & \ensuremath{\textbf{fail}} \\
\rulename{app-tup}  & \ensuremath{\langle{}\Varid{v}_{\mathrm{0}}\;\cdots\;\Varid{v}_{\Varid{n}}\rangle\,\Varid{v}} & \movesto
	& \ensuremath{\exists\Varid{x}.\,\Varid{x}=\Varid{v};\,(\Varid{x}\mathrel{=}\mathrm{0};\,\Varid{v}_{\mathrm{0}}\;\rule[-0.3mm]{0.5mm}{2.5mm}\;\cdots\;\rule[-0.3mm]{0.5mm}{2.5mm}\;\Varid{x}\mathrel{=}\Varid{n};\,\Varid{v}_{\Varid{n}})} & if $\ensuremath{\Varid{x}} \not \in \freevars{\ensuremath{\Varid{v}}}, n \geq 0$ \\
\rulename{app-add}  & \ensuremath{\textbf{add}\langle\Varid{k}_{\mathrm{1}},\Varid{k}_{\mathrm{2}}\rangle} & \movesto & \ensuremath{\Varid{k}_{\mathrm{1}}\mathbin{+}\Varid{k}_{\mathrm{2}}} \\
\rulename{app-gt}      & \ensuremath{\textbf{gt}\langle\Varid{k}_{\mathrm{1}},\Varid{k}_{\mathrm{2}}\rangle}   & \movesto & \ensuremath{\Varid{k}_{\mathrm{1}}} &  \text{if $k_1 > k_2$} \\
\rulename{app-gt-fail} & \ensuremath{\textbf{gt}\langle\Varid{k}_{\mathrm{1}},\Varid{k}_{\mathrm{2}}\rangle}   & \movesto & \ensuremath{\textbf{fail}} &  \text{if $k_1 \leq k_2$}
\end{longtable}
\vspace{-4mm}
\begin{longtable}[l]{@{\quad} l @{\;}  r >{$}c<{$} l l}
\rulegroupname{Speculation:~\rspc}
\rulename{choose}     & \ensuremath{\Conid{SX} [\mskip1.5mu CX [\mskip1.5mu \Varid{e}_{\mathrm{1}}\;\rule[-0.3mm]{0.5mm}{2.5mm}\;\Varid{e}_{\mathrm{2}}\mskip1.5mu]\mskip1.5mu]} & \movesto & \ensuremath{\Conid{SX} [\mskip1.5mu CX [\mskip1.5mu \Varid{e}_{\mathrm{1}}\mskip1.5mu]\;\rule[-0.3mm]{0.5mm}{2.5mm}\;CX [\mskip1.5mu \Varid{e}_{\mathrm{2}}\mskip1.5mu]\mskip1.5mu]} & if $\ensuremath{CX} \ne \hole$ \\
\rulename{choose-assoc} & \ensuremath{\Conid{SX} [\mskip1.5mu (\Varid{e}_{\mathrm{1}}\;\rule[-0.3mm]{0.5mm}{2.5mm}\;\Varid{e}_{\mathrm{2}})\;\rule[-0.3mm]{0.5mm}{2.5mm}\;\Varid{e}_{\mathrm{3}}\mskip1.5mu]} & \movesto & \ensuremath{\Conid{SX} [\mskip1.5mu \Varid{e}_{\mathrm{1}}\;\rule[-0.3mm]{0.5mm}{2.5mm}\;(\Varid{e}_{\mathrm{2}}\;\rule[-0.3mm]{0.5mm}{2.5mm}\;\Varid{e}_{\mathrm{3}})\mskip1.5mu]} \\
\rulename{choose-r}   & \ensuremath{\Conid{SX} [\mskip1.5mu \textbf{fail}\;\rule[-0.3mm]{0.5mm}{2.5mm}\;\Varid{e}\mskip1.5mu]}       & \movesto & \ensuremath{\Conid{SX} [\mskip1.5mu \Varid{e}\mskip1.5mu]} \\
\rulename{choose-l}   & \ensuremath{\Conid{SX} [\mskip1.5mu \Varid{e}\;\rule[-0.3mm]{0.5mm}{2.5mm}\;\textbf{fail}\mskip1.5mu]}       & \movesto & \ensuremath{\Conid{SX} [\mskip1.5mu \Varid{e}\mskip1.5mu]} \\
\rulename{one-fail}   & \ensuremath{\textbf{one}\{\textbf{fail}\}}           & \movesto & \ensuremath{\textbf{fail}} \\
\rulename{one-choice} & \ensuremath{\textbf{one}\{\Varid{e}_{\mathrm{1}}\hspace{1mm}\mathop{\rule[-0.3mm]{0.5mm}{2.5mm}}\hspace{1mm}\Varid{e}_{\mathrm{2}}\}} & \movesto & \ensuremath{\Varid{e}_{\mathrm{1}}} & if \wff{\emptyset}{\ensuremath{\Varid{e}_{\mathrm{1}}}}{\ensuremath{\overline{x}}}{\ensuremath{\overline{c}}}{\ensuremath{\Varid{v}}} \\
\rulename{one-value}  & \ensuremath{\textbf{one}\{\Varid{e}\}}              & \movesto & \ensuremath{\Varid{e}} & if \wff{\emptyset}{\ensuremath{\Varid{e}}}{\ensuremath{\overline{x}}}{\ensuremath{\overline{c}}}{\ensuremath{\Varid{v}}} \\
\rulename{all-fail}   & \ensuremath{\textbf{all}\{\textbf{fail}\}}           & \movesto & \ensuremath{\langle{}\rangle} \\
\rulename{all-choice} & \ensuremath{\textbf{all}\{\Varid{e}_{\mathrm{1}}\;\rule[-0.3mm]{0.5mm}{2.5mm}\;\cdots\;\rule[-0.3mm]{0.5mm}{2.5mm}\;\Varid{e}_{\Varid{n}}\}} & \movesto &
   \ensuremath{\exists\overline{x}.\,\overline{c};\,\langle{}\overline{v}\rangle} &
               if \wffmany{\ensuremath{\overline{e}}}{\ensuremath{\overline{x}}}{\ensuremath{\overline{c}}}{\overline{\ensuremath{\Varid{v}}}}, $n \geq 1$ \\
%% \rulename{all-value}  & |all e|              & \movesto & |def xs cs ; tup v|     & if \wffv{e}{|xs|}{|cs|}{|v|}
\end{longtable}
\caption{\textbf{The Verse Calculus: Rewrite Rules}}
   \label{fig:rewrite-rules} \label{fig:choice-rules} \label{fig:contexts}
\end{figure}
}

\newcommand{\figureNormalizationRules}{
\begin{figure}
\begin{longtable}{@{\quad} l r >{$}c<{$} l l}
\rulegroupname{Normalization:~\rcon}
\rulename{norm-val}       & \ensuremath{\Varid{v};\,\Varid{e}}               & \movesto & \ensuremath{\Varid{e}}                   & \\
\rulename{norm-seq-assoc} & \ensuremath{(eu;\,\Varid{e}_{\mathrm{1}});\,\Varid{e}_{\mathrm{2}}}   & \movesto & \ensuremath{eu;\,(\Varid{e}_{\mathrm{1}};\,\Varid{e}_{\mathrm{2}})}    & \\
\rulename{norm-seq-swap1} & \ensuremath{eu;\,(\Varid{x}\mathrel{=}\Varid{v};\,\Varid{e})}    & \movesto & \ensuremath{\Varid{x}=\Varid{v};\,(eu;\,\Varid{e})}     & if \ensuremath{eu} not of form \ensuremath{\Varid{x'}=\Varid{v'}}\\
\rulename{norm-seq-swap2} & \ensuremath{eu;\,(\Varid{x}\mathrel{=}\Varid{s};\,\Varid{e})}    & \movesto & \ensuremath{\Varid{x}=\Varid{s};\,(eu;\,\Varid{e})}     & if \ensuremath{eu} not of form \ensuremath{\Varid{x'}=\Varid{s'}}\\
% \rulename{norm-eq}        & |e ; x=v|            & \movesto & |e ; (x=v; x)|            & \\
\rulename{norm-eq-swap}   & \ensuremath{\mathit{hnf}=\Varid{x}}            & \movesto & \ensuremath{\Varid{x}=\mathit{hnf}}            & \\
\rulename{norm-seq-defr}  & \ensuremath{(\exists\Varid{x}.\,\Varid{e}_{\mathrm{1}});\,\Varid{e}_{\mathrm{2}}}    & \movesto & \ensuremath{\exists\Varid{x}.\,(\Varid{e}_{\mathrm{1}};\,\Varid{e}_{\mathrm{2}})}    & if $\ensuremath{\Varid{x}} \not \in \freevars{\ensuremath{\Varid{e}_{\mathrm{2}}}}$\\
\rulename{norm-seq-defl}  & \ensuremath{eu;\,(\exists\Varid{x}.\,\Varid{e})}   & \movesto & \ensuremath{\exists\Varid{x}.\,eu;\,\Varid{e}}    & if $\ensuremath{\Varid{x}} \not \in \freevars{\ensuremath{eu}}$\\
\rulename{norm-defr}      & \ensuremath{\Varid{v}=(\exists\Varid{y}.\,\Varid{e}_{\mathrm{1}});\,\Varid{e}_{\mathrm{2}}}      & \movesto & \ensuremath{\exists\Varid{y}.\,\Varid{v}=\Varid{e}_{\mathrm{1}};\,\Varid{e}_{\mathrm{2}}}   & if $\ensuremath{\Varid{y}} \not \in \freevars{\ensuremath{\Varid{v},\Varid{e}_{\mathrm{2}}}}$ \\
\rulename{norm-seqr}      & \ensuremath{\Varid{v}=(eu;\,\Varid{e}_{\mathrm{1}});\,\Varid{e}_{\mathrm{2}}}    & \movesto & \ensuremath{eu;\,\Varid{v}=\Varid{e}_{\mathrm{1}};\,\Varid{e}_{\mathrm{2}}} \\
% \rulename{norm-unifyr}    & |v_1 == (v_2 == e_1); e_2|   & \movesto & |def x (x == v_1; x == v_2; x == e_1; e_2)|   & $|x|$ fresh \\[0.05in]

\rulegroupname{Fail Propagation:~\rfail}
\rulename{fail-seql}      & \ensuremath{\textbf{fail};\,\Varid{e}}            & \movesto & \ensuremath{\textbf{fail}} \\
\rulename{fail-seqr}      & \ensuremath{\Varid{e};\,\textbf{fail}}            & \movesto & \ensuremath{\textbf{fail}} \\
\rulename{fail-eq}        & \ensuremath{\Varid{v}=\textbf{fail}}           & \movesto & \ensuremath{\textbf{fail}} \\[0.5mm]
\rulegroupname{Garbage Collection:~$\rgc$}
\rulename{elim-def}     & \ensuremath{\Varid{e}_{\mathrm{1}}} & \movesto & \multicolumn{2}{l}
         {\ensuremath{\Varid{e}_{\mathrm{2}}} \quad \quad if $\wff{\emptyset}{\ensuremath{\Varid{e}_{\mathrm{1}}}}{\ensuremath{\overline{x}}}{\ensuremath{\overline{c}}}{\ensuremath{\Varid{e}_{\mathrm{2}}}}$
             and $\ensuremath{\overline{x}} \not\in \freevars{\ensuremath{\Varid{e}_{\mathrm{2}}}}$} \\
\rulegroupname{Structural rules}
        \rulename{swap-d} & \ensuremath{\exists\Varid{x}.\,\exists\Varid{y}.\,\Varid{e}} & \equiv & \ensuremath{\exists\Varid{y}.\,\exists\Varid{x}.\,\Varid{e}} &  \\
        \rulename{swap-c} & \ensuremath{\Varid{x}_{\mathrm{1}}\mathrel{=}\Varid{v}_{\mathrm{1}};\,\Varid{x}_{\mathrm{2}}\mathrel{=}\Varid{v}_{\mathrm{2}};\,\Varid{e}} & \equiv & \ensuremath{\Varid{x}_{\mathrm{2}}\mathrel{=}\Varid{v}_{\mathrm{2}};\,\Varid{x}_{\mathrm{1}}\mathrel{=}\Varid{v}_{\mathrm{1}};\,\Varid{e}}
\end{longtable}
\caption{\textbf{The Verse Calculus: Normalization Rules}}\label{fig:normalization-rules}

\end{figure}
%
\begin{figure}
% \judgementHead{Wellformedness of Results}{\wf{r}}

\begin{mathpar}
\fbox{\wff{\Gamma}{\ensuremath{\Varid{e}_{\mathrm{1}}}}{\ensuremath{\overline{x}}}{\ensuremath{\overline{c}}}{\ensuremath{\Varid{e}_{\mathrm{2}}}}}

\inferrule
  { }
  { \wff{\Gamma}{\ensuremath{\Varid{e}}}{\emptyset}{\emptyset}{\ensuremath{\Varid{e}}} }
  { \rulename{WF-Exp}}

% \inferrule
%   { x \in \Gamma }
%  {\wff{\Gamma}{|x == v|}{\emptyset}{|x == v|}{|x|}}
%  {\rulename{WF-Eq-Val}}
\\
\inferrule
  { \wff{\Gamma, x}{e_1}{\ensuremath{\overline{x}}}{\ensuremath{\overline{c}}}{\ensuremath{\Varid{e}_{\mathrm{2}}}} \\
    x \not \in \ensuremath{\overline{x}}}
  { \wff{\Gamma}{\ensuremath{\exists\Varid{x}.\,\Varid{e}_{\mathrm{1}}}}{\ensuremath{\Varid{x},\overline{x}}}{\ensuremath{\overline{c}}}{\ensuremath{\Varid{e}_{\mathrm{2}}}} }
  { \rulename{WF-Def} }

\inferrule
  { x \in \Gamma \quad v \not= x \\
    \text{if $v = s$ then $x \not\in \freevars{\ensuremath{\Varid{e}_{\mathrm{1}}}}$} \\\\
    \wff{\Gamma - x}{\ensuremath{\Varid{e}_{\mathrm{1}}}}{\ensuremath{\overline{x}}}{\ensuremath{\overline{c}}}{\ensuremath{\Varid{e}_{\mathrm{2}}}} \\
    \freevars{h} \not\in \ensuremath{\overline{x}}
  }
  {\wff{\Gamma}{\ensuremath{\Varid{x}=\Varid{v};\,\Varid{e}_{\mathrm{1}}}}{\ensuremath{\overline{x}}}{\ensuremath{\Varid{x}=\Varid{v},\overline{c}}}{\ensuremath{\Varid{e}_{\mathrm{2}}}}}
  {\rulename{WF-Eq}}

\\

%\inferrule
%  { \wff{\Gamma}{e_1}{|xs|}{|cs|}{|e_2|} \quad
%    \wff{\Gamma,|xs|}{|es_1|}{|xs|'}{|cs|'}{|es_2|}
%  }
%  { \wff{\Gamma}{|e_1|,|es_1|}{|xs|,|xs|'}{|cs|,|cs|'}{|e_2,es_2|} }
%  {\rulename{WF-many}}
\inferrule
  { \wff{\emptyset}{r_1}{\ensuremath{\overline{x}}_1}{\ensuremath{\overline{c}}_1}{\ensuremath{\Varid{e}_{\mathrm{1}}}} \quad
    \ensuremath{\cdots} \quad
    \wff{\emptyset}{r_n}{\ensuremath{\overline{x}}_n}{\ensuremath{\overline{c}}_n}{\ensuremath{\Varid{e}_{\Varid{n}}}} \quad
    \quad
    \text{all \ensuremath{\Varid{x}_{\Varid{i}}} distinct}
  }
  { \wffmany{\ensuremath{\Varid{r}_{\mathrm{1}},\cdots,\Varid{r}_{\Varid{n}}}}{\ensuremath{\overline{x}}_1, \ensuremath{\cdots,\overline{x}}_n}{\ensuremath{\overline{c}}_1, \ensuremath{\cdots,\overline{c}}_n}{\ensuremath{\Varid{e}_{\mathrm{1}},\cdots,\Varid{e}_{\Varid{n}}}} }
  {\rulename{WF-many} }
\end{mathpar}
\caption{\textbf{Well-formedness of Results}}\label{fig:wf}
\end{figure}
}
